\IfFileExists{stacks-project.cls}{%
\documentclass{stacks-project}
}{%
\documentclass{amsart}
}

% For dealing with references we use the comment environment
\usepackage{verbatim}
\newenvironment{reference}{\comment}{\endcomment}
%\newenvironment{reference}{}{}
\newenvironment{slogan}{\comment}{\endcomment}
\newenvironment{history}{\comment}{\endcomment}

% For commutative diagrams we use Xy-pic
\usepackage[all]{xy}

% We use 2cell for 2-commutative diagrams.
\xyoption{2cell}
\UseAllTwocells

% We use multicol for the list of chapters between chapters
\usepackage{multicol}

% This is generally recommended for better output
\usepackage{lmodern}
\usepackage[T1]{fontenc}

% For cross-file-references
\usepackage{xr-hyper}

% Package for hypertext links:
\usepackage{hyperref}

% For any local file, say "hello.tex" you want to link to please
% use \externaldocument[hello-]{hello}
\externaldocument[introduction-]{introduction}
\externaldocument[conventions-]{conventions}
\externaldocument[sets-]{sets}
\externaldocument[categories-]{categories}
\externaldocument[topology-]{topology}
\externaldocument[sheaves-]{sheaves}
\externaldocument[sites-]{sites}
\externaldocument[stacks-]{stacks}
\externaldocument[fields-]{fields}
\externaldocument[algebra-]{algebra}
\externaldocument[brauer-]{brauer}
\externaldocument[homology-]{homology}
\externaldocument[derived-]{derived}
\externaldocument[simplicial-]{simplicial}
\externaldocument[more-algebra-]{more-algebra}
\externaldocument[smoothing-]{smoothing}
\externaldocument[modules-]{modules}
\externaldocument[sites-modules-]{sites-modules}
\externaldocument[injectives-]{injectives}
\externaldocument[cohomology-]{cohomology}
\externaldocument[sites-cohomology-]{sites-cohomology}
\externaldocument[dga-]{dga}
\externaldocument[dpa-]{dpa}
\externaldocument[sdga-]{sdga}
\externaldocument[hypercovering-]{hypercovering}
\externaldocument[schemes-]{schemes}
\externaldocument[constructions-]{constructions}
\externaldocument[properties-]{properties}
\externaldocument[morphisms-]{morphisms}
\externaldocument[coherent-]{coherent}
\externaldocument[divisors-]{divisors}
\externaldocument[limits-]{limits}
\externaldocument[varieties-]{varieties}
\externaldocument[topologies-]{topologies}
\externaldocument[descent-]{descent}
\externaldocument[perfect-]{perfect}
\externaldocument[more-morphisms-]{more-morphisms}
\externaldocument[flat-]{flat}
\externaldocument[groupoids-]{groupoids}
\externaldocument[more-groupoids-]{more-groupoids}
\externaldocument[etale-]{etale}
\externaldocument[chow-]{chow}
\externaldocument[intersection-]{intersection}
\externaldocument[pic-]{pic}
\externaldocument[weil-]{weil}
\externaldocument[adequate-]{adequate}
\externaldocument[dualizing-]{dualizing}
\externaldocument[duality-]{duality}
\externaldocument[discriminant-]{discriminant}
\externaldocument[derham-]{derham}
\externaldocument[local-cohomology-]{local-cohomology}
\externaldocument[algebraization-]{algebraization}
\externaldocument[curves-]{curves}
\externaldocument[resolve-]{resolve}
\externaldocument[models-]{models}
\externaldocument[functors-]{functors}
\externaldocument[equiv-]{equiv}
\externaldocument[pione-]{pione}
\externaldocument[etale-cohomology-]{etale-cohomology}
\externaldocument[proetale-]{proetale}
\externaldocument[relative-cycles-]{relative-cycles}
\externaldocument[more-etale-]{more-etale}
\externaldocument[trace-]{trace}
\externaldocument[crystalline-]{crystalline}
\externaldocument[spaces-]{spaces}
\externaldocument[spaces-properties-]{spaces-properties}
\externaldocument[spaces-morphisms-]{spaces-morphisms}
\externaldocument[decent-spaces-]{decent-spaces}
\externaldocument[spaces-cohomology-]{spaces-cohomology}
\externaldocument[spaces-limits-]{spaces-limits}
\externaldocument[spaces-divisors-]{spaces-divisors}
\externaldocument[spaces-over-fields-]{spaces-over-fields}
\externaldocument[spaces-topologies-]{spaces-topologies}
\externaldocument[spaces-descent-]{spaces-descent}
\externaldocument[spaces-perfect-]{spaces-perfect}
\externaldocument[spaces-more-morphisms-]{spaces-more-morphisms}
\externaldocument[spaces-flat-]{spaces-flat}
\externaldocument[spaces-groupoids-]{spaces-groupoids}
\externaldocument[spaces-more-groupoids-]{spaces-more-groupoids}
\externaldocument[bootstrap-]{bootstrap}
\externaldocument[spaces-pushouts-]{spaces-pushouts}
\externaldocument[spaces-chow-]{spaces-chow}
\externaldocument[groupoids-quotients-]{groupoids-quotients}
\externaldocument[spaces-more-cohomology-]{spaces-more-cohomology}
\externaldocument[spaces-simplicial-]{spaces-simplicial}
\externaldocument[spaces-duality-]{spaces-duality}
\externaldocument[formal-spaces-]{formal-spaces}
\externaldocument[restricted-]{restricted}
\externaldocument[spaces-resolve-]{spaces-resolve}
\externaldocument[formal-defos-]{formal-defos}
\externaldocument[defos-]{defos}
\externaldocument[cotangent-]{cotangent}
\externaldocument[examples-defos-]{examples-defos}
\externaldocument[algebraic-]{algebraic}
\externaldocument[examples-stacks-]{examples-stacks}
\externaldocument[stacks-sheaves-]{stacks-sheaves}
\externaldocument[criteria-]{criteria}
\externaldocument[artin-]{artin}
\externaldocument[quot-]{quot}
\externaldocument[stacks-properties-]{stacks-properties}
\externaldocument[stacks-morphisms-]{stacks-morphisms}
\externaldocument[stacks-limits-]{stacks-limits}
\externaldocument[stacks-cohomology-]{stacks-cohomology}
\externaldocument[stacks-perfect-]{stacks-perfect}
\externaldocument[stacks-introduction-]{stacks-introduction}
\externaldocument[stacks-more-morphisms-]{stacks-more-morphisms}
\externaldocument[stacks-geometry-]{stacks-geometry}
\externaldocument[moduli-]{moduli}
\externaldocument[moduli-curves-]{moduli-curves}
\externaldocument[examples-]{examples}
\externaldocument[exercises-]{exercises}
\externaldocument[guide-]{guide}
\externaldocument[desirables-]{desirables}
\externaldocument[coding-]{coding}
\externaldocument[obsolete-]{obsolete}
\externaldocument[fdl-]{fdl}
\externaldocument[index-]{index}

% Theorem environments.
%
\theoremstyle{plain}
\newtheorem{theorem}[subsection]{Theorem}
\newtheorem{proposition}[subsection]{Proposition}
\newtheorem{lemma}[subsection]{Lemma}

\theoremstyle{definition}
\newtheorem{definition}[subsection]{Definition}
\newtheorem{example}[subsection]{Example}
\newtheorem{exercise}[subsection]{Exercise}
\newtheorem{situation}[subsection]{Situation}

\theoremstyle{remark}
\newtheorem{remark}[subsection]{Remark}
\newtheorem{remarks}[subsection]{Remarks}

\numberwithin{equation}{subsection}

% Macros
%
\def\lim{\mathop{\mathrm{lim}}\nolimits}
\def\colim{\mathop{\mathrm{colim}}\nolimits}
\def\Spec{\mathop{\mathrm{Spec}}}
\def\Hom{\mathop{\mathrm{Hom}}\nolimits}
\def\Ext{\mathop{\mathrm{Ext}}\nolimits}
\def\SheafHom{\mathop{\mathcal{H}\!\mathit{om}}\nolimits}
\def\SheafExt{\mathop{\mathcal{E}\!\mathit{xt}}\nolimits}
\def\Sch{\mathit{Sch}}
\def\Mor{\mathop{\mathrm{Mor}}\nolimits}
\def\Ob{\mathop{\mathrm{Ob}}\nolimits}
\def\Sh{\mathop{\mathit{Sh}}\nolimits}
\def\NL{\mathop{N\!L}\nolimits}
\def\CH{\mathop{\mathrm{CH}}\nolimits}
\def\proetale{{pro\text{-}\acute{e}tale}}
\def\etale{{\acute{e}tale}}
\def\QCoh{\mathit{QCoh}}
\def\Ker{\mathop{\mathrm{Ker}}}
\def\Im{\mathop{\mathrm{Im}}}
\def\Coker{\mathop{\mathrm{Coker}}}
\def\Coim{\mathop{\mathrm{Coim}}}

% Boxtimes
%
\DeclareMathSymbol{\boxtimes}{\mathbin}{AMSa}{"02}

%
% Macros for moduli stacks/spaces
%
\def\QCohstack{\mathcal{QC}\!\mathit{oh}}
\def\Cohstack{\mathcal{C}\!\mathit{oh}}
\def\Spacesstack{\mathcal{S}\!\mathit{paces}}
\def\Quotfunctor{\mathrm{Quot}}
\def\Hilbfunctor{\mathrm{Hilb}}
\def\Curvesstack{\mathcal{C}\!\mathit{urves}}
\def\Polarizedstack{\mathcal{P}\!\mathit{olarized}}
\def\Complexesstack{\mathcal{C}\!\mathit{omplexes}}
% \Pic is the operator that assigns to X its picard group, usage \Pic(X)
% \Picardstack_{X/B} denotes the Picard stack of X over B
% \Picardfunctor_{X/B} denotes the Picard functor of X over B
\def\Pic{\mathop{\mathrm{Pic}}\nolimits}
\def\Picardstack{\mathcal{P}\!\mathit{ic}}
\def\Picardfunctor{\mathrm{Pic}}
\def\Deformationcategory{\mathcal{D}\!\mathit{ef}}
\usepackage{cleveref}
\crefname{originallemma}{original lemma}{original lemmas}
\Crefname{originallemma}{Original lemma}{Original lemmas}
\makeatletter
\newlabel{lemma-largest-open-of-the-source}{{0CEY}{source}{}{sourcepointer-1}{}}
\newlabel{lemma-largest-open-of-the-source@cref}{{[originallemma][0][]0CEY}{[1][0][]source}{}{}{}}
\makeatother
\crefname{originallemma}{original lemma}{original lemmas}
\Crefname{originallemma}{Original lemma}{Original lemmas}
\makeatletter
\newlabel{lemma-local-source-target-implies}{{04R0}{source}{}{sourcepointer-2}{}}
\newlabel{lemma-local-source-target-implies@cref}{{[originallemma][0][]04R0}{[1][0][]source}{}{}{}}
\makeatother
\title{Etale localness on the source-and-target is compatible with etale base change of the diagram}
\author{The Stacks Project Authors}
\date{stacks-project:descent.tex:local-snapshot-20261012; GNU FDL 1.2 or later}
\begin{document}
\maketitle
\noindent\textit{Editorial note.} All delivered mathematical text is the author's own, copied from the pinned Stacks Project chapter file descent.tex; no statement or proof step was rewritten, reordered or removed.\par\medskip
\noindent\textit{Editorial note.} The counted spans carry 3 navigation stubs for the named definition and lemmas of the localness notion; the application of each clause is written out inside the span, so nothing load bearing is delegated.\par\medskip
\noindent\textit{Editorial note.} The preamble span points at the chapter's own preamble.tex shipped with the frozen source, used verbatim as staged.\par\medskip
\noindent\textit{Editorial note.} The label lemma-etale-etale-local-source-target is defined only in descent.tex among the staged chapters, so no same-name restatement had to be weighed for this card.\par\medskip
\noindent\textit{Allowed standard prerequisites (editorial).}\begin{itemize}
\item Etale localness on the source-and-target and its equivalent clauses. Same chapter, Definition definition-local-source-target and Lemma lemma-local-source-target-implies, invoked at proof lines 7812-7825 and 7834-7842. Supplies the clauses (1), (3) and (4) used in both directions.
\item The locus W(f) attached to a property via the largest open on which the property holds. Same chapter, Lemma lemma-largest-open-of-the-source, invoked at proof lines 7808 and 7812-7813. Supplies the object being compared and the input to which the localness definition is applied.
\item An etale morphism is open. Morphisms, the openness of etale morphisms, used at proof lines 7838-7840. Supplies that the image of the etale neighbourhood is open, so that a single-member etale covering is obtained.
\end{itemize}\medskip
\ifdefined\hypertarget\hypertarget{sourcepointer-1}{}\fi
\noindent\textit{Source reference.} Original lemma 0CEY: sources/stacks-descent/pinned-original-source.tex, source lines 6890--6890. Pointer outside the counted span, invoked once for context; the load-bearing argument of the counted statement is the author's proof reproduced in full in the proof span.\par\medskip
\ifdefined\hypertarget\hypertarget{sourcepointer-2}{}\fi
\noindent\textit{Source reference.} Original lemma 04R0: sources/stacks-descent/pinned-original-source.tex, source lines 7399--7399. Pointer outside the counted span, invoked once for context; the load-bearing argument of the counted statement is the author's proof reproduced in full in the proof span.\par\medskip
\begin{definition}
\label{definition-local-source-target}
Let $\mathcal{P}$ be a property of morphisms of schemes.
We say $\mathcal{P}$ is {\it \'etale local on source-and-target} if
\begin{enumerate}
\item (stable under precomposing with \'etale maps)
if $f : X \to Y$ is \'etale and $g : Y \to Z$ has $\mathcal{P}$,
then $g \circ f$ has $\mathcal{P}$,
\item (stable under \'etale base change)
if $f : X \to Y$ has $\mathcal{P}$ and $Y' \to Y$ is \'etale, then
the base change $f' : Y' \times_Y X \to Y'$ has $\mathcal{P}$, and
\item (locality) given a morphism $f : X \to Y$ the following are equivalent
\begin{enumerate}
\item $f$ has $\mathcal{P}$,
\item for every $x \in X$ there exists a commutative diagram
$$
\xymatrix{
U \ar[d]_a \ar[r]_h & V \ar[d]^b \\
X \ar[r]^f & Y
}
$$
with \'etale vertical arrows and $u \in U$ with $a(u) = x$ such that
$h$ has $\mathcal{P}$.
\end{enumerate}
\end{enumerate}
\end{definition}

\begin{lemma}
\label{lemma-etale-etale-local-source-target}
Let $\mathcal{P}$ be a property of morphisms of schemes which
is \'etale local on the source-and-target.
Given a commutative diagram of schemes
$$
\vcenter{
\xymatrix{
X' \ar[d]_{g'} \ar[r]_{f'} & Y' \ar[d]^g \\
X \ar[r]^f & Y
}
}
\quad\text{with points}\quad
\vcenter{
\xymatrix{
x' \ar[d] \ar[r] & y' \ar[d] \\
x \ar[r] & y
}
}
$$
such that $g'$ is \'etale at $x'$ and $g$ is \'etale at $y'$, then
$x \in W(f) \Leftrightarrow x' \in W(f')$
where $W(-)$ is as in Lemma \ref{lemma-largest-open-of-the-source}.
\end{lemma}

\begin{proof}
Lemma \ref{lemma-largest-open-of-the-source} applies since
$\mathcal{P}$ is \'etale local on the source by
Lemma \ref{lemma-local-source-target-implies}.

\medskip\noindent
Assume $x \in W(f)$. Let $U' \subset X'$ and $V' \subset Y'$
be open neighbourhoods of $x'$ and $y'$ such that $f'(U') \subset V'$,
$g'(U') \subset W(f)$ and $g'|_{U'}$ and $g|_{V'}$ are \'etale.
Then $f \circ g'|_{U'} = g \circ f'|_{U'}$
has $\mathcal{P}$ by property (1) of
Definition \ref{definition-local-source-target}.
Then $f'|_{U'} : U' \to V'$ has property $\mathcal{P}$
by (4) of Lemma \ref{lemma-local-source-target-implies}.
Then by (3) of Lemma \ref{lemma-local-source-target-implies}
we conclude that $f'_{U'} : U' \to Y'$ has $\mathcal{P}$.
Hence $U' \subset W(f')$ by definition. Hence $x' \in W(f')$.

\medskip\noindent
Assume $x' \in W(f')$. Let $U' \subset X'$ and $V' \subset Y'$
be open neighbourhoods of $x'$ and $y'$ such that $f'(U') \subset V'$,
$U' \subset W(f')$ and $g'|_{U'}$ and $g|_{V'}$ are \'etale.
Then $U' \to Y'$ has $\mathcal{P}$ by definition of $W(f')$.
Then $U' \to V'$ has $\mathcal{P}$ by (4) of
Lemma \ref{lemma-local-source-target-implies}.
Then $U' \to Y$ has $\mathcal{P}$ by (3) of
Lemma \ref{lemma-local-source-target-implies}.
Let $U \subset X$ be the image of the \'etale (hence open)
morphism $g'|_U' : U' \to X$. Then $\{U' \to U\}$
is an \'etale covering and we conclude that
$U \to Y$ has $\mathcal{P}$ by (1) of
Lemma \ref{lemma-local-source-target-implies}.
Thus $U \subset W(f)$ by definition. Hence $x \in W(f)$.
\end{proof}

\end{document}
